\documentclass[10pt,a4paper]{amsart}
\usepackage{amsmath,amssymb,amsthm,graphicx}
\usepackage{verbatim}
\newenvironment{reference}{\comment}{\endcomment}
%\newenvironment{reference}{}{}
\newenvironment{slogan}{\comment}{\endcomment}
\newenvironment{history}{\comment}{\endcomment}
% For commutative diagrams we use Xy-pic
\usepackage[all]{xy}

% We use 2cell for 2-commutative diagrams.
\xyoption{2cell}
% Theorem environments.
%
\theoremstyle{plain}
\newtheorem{theorem}[subsection]{Theorem}
\newtheorem{proposition}[subsection]{Proposition}
\newtheorem{lemma}[subsection]{Lemma}

\theoremstyle{definition}
\newtheorem{definition}[subsection]{Definition}
\newtheorem{example}[subsection]{Example}
\newtheorem{exercise}[subsection]{Exercise}
\newtheorem{situation}[subsection]{Situation}

\theoremstyle{remark}
\newtheorem{remark}[subsection]{Remark}
\newtheorem{remarks}[subsection]{Remarks}

\numberwithin{equation}{subsection}

% Macros
%
\def\lim{\mathop{\mathrm{lim}}\nolimits}
\def\colim{\mathop{\mathrm{colim}}\nolimits}
\def\Spec{\mathop{\mathrm{Spec}}}
\def\Hom{\mathop{\mathrm{Hom}}\nolimits}
\def\Ext{\mathop{\mathrm{Ext}}\nolimits}
\def\SheafHom{\mathop{\mathcal{H}\!\mathit{om}}\nolimits}
\def\SheafExt{\mathop{\mathcal{E}\!\mathit{xt}}\nolimits}
\def\Sch{\mathit{Sch}}
\def\Mor{\mathop{\mathrm{Mor}}\nolimits}
\def\Ob{\mathop{\mathrm{Ob}}\nolimits}
\def\Sh{\mathop{\mathit{Sh}}\nolimits}
\def\NL{\mathop{N\!L}\nolimits}
\def\CH{\mathop{\mathrm{CH}}\nolimits}
\def\proetale{{pro\text{-}\acute{e}tale}}
\def\etale{{\acute{e}tale}}
\def\QCoh{\mathit{QCoh}}
\def\Ker{\mathop{\mathrm{Ker}}}
\def\Im{\mathop{\mathrm{Im}}}
\def\Coker{\mathop{\mathrm{Coker}}}
\def\Coim{\mathop{\mathrm{Coim}}}

% Boxtimes
%
\DeclareMathSymbol{\boxtimes}{\mathbin}{AMSa}{"02}

%
% Macros for moduli stacks/spaces
%
\def\QCohstack{\mathcal{QC}\!\mathit{oh}}
\def\Cohstack{\mathcal{C}\!\mathit{oh}}
\def\Spacesstack{\mathcal{S}\!\mathit{paces}}
\def\Quotfunctor{\mathrm{Quot}}
\def\Hilbfunctor{\mathrm{Hilb}}
\def\Curvesstack{\mathcal{C}\!\mathit{urves}}
\def\Polarizedstack{\mathcal{P}\!\mathit{olarized}}
\def\Complexesstack{\mathcal{C}\!\mathit{omplexes}}
% \Pic is the operator that assigns to X its picard group, usage \Pic(X)
% \Picardstack_{X/B} denotes the Picard stack of X over B
% \Picardfunctor_{X/B} denotes the Picard functor of X over B
\def\Pic{\mathop{\mathrm{Pic}}\nolimits}
\def\Picardstack{\mathcal{P}\!\mathit{ic}}
\def\Picardfunctor{\mathrm{Pic}}
\def\Deformationcategory{\mathcal{D}\!\mathit{ef}}
\usepackage{cleveref}
\crefname{originallemma}{original lemma}{original lemmas}
\Crefname{originallemma}{Original lemma}{Original lemmas}
\makeatletter
\newlabel{lemma-closed-in-quasi-compact}{{lemma-closed-in-quasi-compact}{source}{}{sourcepointer-1}{}}
\newlabel{lemma-closed-in-quasi-compact@cref}{{[originallemma][0][]lemma-closed-in-quasi-compact}{[1][0][]source}{}{}{}}
\makeatother
\crefname{originallemma}{original lemma}{original lemmas}
\Crefname{originallemma}{Original lemma}{Original lemmas}
\makeatletter
\newlabel{lemma-constructible-hausdorff-quasi-compact}{{lemma-constructible-hausdorff-quasi-compact}{source}{}{sourcepointer-2}{}}
\newlabel{lemma-constructible-hausdorff-quasi-compact@cref}{{[originallemma][0][]lemma-constructible-hausdorff-quasi-compact}{[1][0][]source}{}{}{}}
\makeatother
\crefname{originallemma}{original lemma}{original lemmas}
\Crefname{originallemma}{Original lemma}{Original lemmas}
\makeatletter
\newlabel{lemma-nonempty-limit}{{lemma-nonempty-limit}{source}{}{sourcepointer-3}{}}
\newlabel{lemma-nonempty-limit@cref}{{[originallemma][0][]lemma-nonempty-limit}{[1][0][]source}{}{}{}}
\makeatother
\title{A cofiltered limit of nonempty spectral spaces with closed transition images is nonempty}
\author{The Stacks Project Authors}
\date{stacks-project:topology.tex:local-snapshot-20261012; GNU FDL 1.2 or later}
\begin{document}
\maketitle
\noindent\textit{Editorial note.} All delivered mathematical text is the author's own, copied from the pinned Stacks Project chapter source sources/stacks-topology/pinned-original-source.tex; no statement or proof step was rewritten, reordered or removed.\par\medskip
\noindent\textit{Editorial note.} Preamble provenance: the Stacks Project keeps one single shared preamble for every chapter; this card takes it directly from the staged sources/stacks-topology/preamble.tex as three verbatim spans (lines 8-12, 14-18 and 154-217), so no derived file was written and nothing was added or redefined. The wrapper class and the base packages are supplied by the assembler.\par\medskip
\noindent\textit{Editorial note.} Bibliography provenance: neither the counted statement nor its proof contains a citation (this lane rescans the whole span before writing the card), so the card carries no bibliography and no entry was reconstructed.\par\medskip
\noindent\textit{Editorial note.} Selection provenance: this card was opened under the file-based control contract (CONTROL.md), with the mechanical pre-card gate run on its statement and proof interval before the card was written.\par\medskip
\noindent\textit{Allowed standard prerequisites (editorial).}\begin{itemize}
\item The statement that the limit of nonempty quasi-compact Hausdorff spaces over a cofiltered category is nonempty. Invoked inside the counted proof (source lines 5045-5047). Provides nonemptiness after passing to the constructible topologies.
\item The statement that the constructible topology of a spectral space is quasi-compact Hausdorff, and that spectral maps are continuous for it. Used inside the counted proof (source lines 5040-5044). Puts the inverse system in the setting of the previous nonemptiness result.
\end{itemize}\medskip
\ifdefined\hypertarget\hypertarget{sourcepointer-1}{}\fi
\noindent\textit{Source reference.} Original lemma lemma-closed-in-quasi-compact: sources/stacks-topology/pinned-original-source.tex, source lines 1879--1879. Navigation pointer only: the counted text refers to the Stacks item lemma-closed-in-quasi-compact, which is an existing labeled result of this same chapter and is not part of this card. The pointer carries the target's own label as its number and its definition line as its span; its content is not reproduced here.\par\medskip
\ifdefined\hypertarget\hypertarget{sourcepointer-2}{}\fi
\noindent\textit{Source reference.} Original lemma lemma-constructible-hausdorff-quasi-compact: sources/stacks-topology/pinned-original-source.tex, source lines 4522--4522. Navigation pointer only: the counted text refers to the Stacks item lemma-constructible-hausdorff-quasi-compact, which is an existing labeled result of this same chapter and is not part of this card. The pointer carries the target's own label as its number and its definition line as its span; its content is not reproduced here.\par\medskip
\ifdefined\hypertarget\hypertarget{sourcepointer-3}{}\fi
\noindent\textit{Source reference.} Original lemma lemma-nonempty-limit: sources/stacks-topology/pinned-original-source.tex, source lines 2724--2724. Navigation pointer only: the counted text refers to the Stacks item lemma-nonempty-limit, which is an existing labeled result of this same chapter and is not part of this card. The pointer carries the target's own label as its number and its definition line as its span; its content is not reproduced here.\par\medskip
\begin{lemma}
\label{lemma-inverse-limit-spectral-spaces-nonempty}
Let $\mathcal{I}$ be a cofiltered category. Let $i \mapsto X_i$ be a diagram
of spectral spaces such that for $a : j \to i$ in $\mathcal{I}$
the corresponding map $f_a : X_j \to X_i$ is spectral.
\begin{enumerate}
\item Given nonempty subsets $Z_i \subset X_i$ closed in the constructible
topology with $f_a(Z_j) \subset Z_i$ for all $a : j \to i$ in $\mathcal{I}$,
then $\lim Z_i$ is nonempty.
\item If each $X_i$ is nonempty, then $X = \lim X_i$ is nonempty.
\end{enumerate}
\end{lemma}

\begin{proof}
Denote $X'_i$ the space $X_i$ endowed with the constructible topology
and $Z'_i$ the corresponding subspace of $X'_i$.
Let $a : j \to i$ in $\mathcal{I}$ be a morphism. As $f_a$ is spectral
it defines a continuous map $f_a : X'_j \to X'_i$. Thus
$f_a|_{Z_j} : Z'_j \to Z'_i$ is a continuous map of quasi-compact
Hausdorff spaces (by
Lemmas \ref{lemma-constructible-hausdorff-quasi-compact} and
\ref{lemma-closed-in-quasi-compact}). By
Lemma \ref{lemma-nonempty-limit} the space $\lim Z'_i$ is nonempty.
Since $\lim Z'_i = \lim Z_i$ as sets we conclude.
\end{proof}

\end{document}
